\section{The precise result and its scope}

The usual Desargues theorem says that two nondegenerate triangles are
perspective from a point if and only if they are perspective from a line.
Pascal says that the three intersections of opposite sides of a hexagon
inscribed in a conic are collinear. Homogeneous coordinates include points
at infinity, so parallel affine lines require no exceptional case.
The statements in the question provide the geometric motivation
\cite{desargues,pascal}; all proofs below are supplied explicitly.

Three levels of statement must be distinguished in a first-order arithmetic
treatment.
\begin{enumerate}
\item \textbf{Internal arithmetic theorems.} For arbitrary integer-coded
coordinates and coefficients, iPRA proves the incidence implications and
the correctness of the displayed line and center algorithms. This is
the entire rational projective plane $\mathbb P^2(\Q)$, rather than a
finite numerical sample.
\item \textbf{Universal algebraic certificates.} The two main identities
hold in every commutative ring. Their geometric consequences hold over
commutative fields with the requisite nonzero, or constructive apartness,
witnesses. The proofs use no order or analytic structure.
\item \textbf{A real-coordinate arithmetic schema.}
Section~\ref{sec:real} states the same conclusions for each fixed finite
family of primitive-recursive rational approximation functions satisfying
explicit arithmetic regularity and apartness premises. This is an iPRA
schema, not quantification over a completed set of all real numbers.
\end{enumerate}
No assertion is made that Desargues follows from the incidence axioms of
an arbitrary abstract projective plane: coordinates are part of the
interpretation. Nor is an ambient Lean theorem about a ring automatically
identified with a derivation in the object theory iPRA.

Here are the principal statements; their notation and witness algorithms
are developed below. A triple is \emph{nonzero} when at least one of its
integer components is nonzero.

\begin{theorem}[Internal Desargues, in both directions]\label{thm:des-summary}
Let $a,b,c,A,B,C$ be integer triples with
$\br a b c\ne0$ and $\br A B C\ne0$. Put
\[
 j_1=a\times A,\quad j_2=b\times B,\quad j_3=c\times C,
\]
\[
 X=(a\times b)\times(A\times B),\quad
 Y=(b\times c)\times(B\times C),\quad
 Z=(c\times a)\times(C\times A).
\]
Suppose all $j_i$ and $X,Y,Z$ are nonzero. Then iPRA proves
\begin{align*}
 &\exists o\,[\nz(o)\land o\cdot j_1=o\cdot j_2=o\cdot j_3=0]\\
 &\qquad\longleftrightarrow
 \exists h\,[\nz(h)\land h\cdot X=h\cdot Y=h\cdot Z=0].
\end{align*}
In the forward direction choose $h=\axis(X,Y,Z)$; in the reverse
direction choose $o=\axis(j_1,j_2,j_3)$. The forward implication does not
need the two triangle determinant assumptions.
\end{theorem}

\begin{theorem}[Internal Pascal]\label{thm:pas-summary}
Let $P_1,\ldots,P_6$ be nonzero integer triples and let
\[
 Q(x,y,z)=q_1x^2+q_2y^2+q_3z^2+q_4xy+q_5xz+q_6yz
\]
have an integer coefficient vector $q\ne0$. Suppose $Q(P_i)=0$ for every
$i$. Define
\[
 X=(P_1\times P_2)\times(P_4\times P_5),\quad
 Y=(P_2\times P_3)\times(P_5\times P_6),
\]
\[
 Z=(P_3\times P_4)\times(P_6\times P_1).
\]
Then iPRA proves $\br X Y Z=0$. If $X,Y,Z$ are nonzero, it additionally
proves that $h=\axis(X,Y,Z)$ is nonzero and satisfies
$h\cdot X=h\cdot Y=h\cdot Z=0$.
\end{theorem}

The nonzero double cross products ensure that the side lines and their
specified intersections really exist. The determinant conclusion remains
true even when such objects degenerate, but the zero vector is never
treated as a projective point. Coincident intersections are allowed; a
common line is then constructed, although it need not be unique.
Replacing a repeated vertex by a tangent is a different construction and
is not silently included in the cross-product convention.

